\lecture{14}{Основы теории устойчивости}

\sourcevideo
    {https://www.youtube.com/playlist?list=PLOzRYVm0a65fhKDS8cYIC7YCuVGJ4FOJx}
    {Week 4: Lecture 14: Stability Theory}

Оптимизационный алгоритм можно рассматривать как динамическую
систему, а его сходимость --- как устойчивость соответствующего
равновесия. Сначала это иллюстрируется сравнением градиентного спуска
и метода тяжёлого шара для плоской функции, затем вводятся основные
понятия устойчивости и прямой метод Ляпунова.

\section{Численный пример и метод тяжёлого шара}

Рассмотрим функцию
\[
    f(x)=x^4,
    \qquad
    f'(x)=4x^3,
    \qquad
    f''(x)=12x^2.
\]
Она строго выпукла и имеет единственный минимум
\[
    x^\star=0,
    \qquad
    f^\star=0,
\]
но не является глобально сильно выпуклой, поскольку
\[
    \inf_{x\in\RR}f''(x)=0,
\]
и не имеет глобально липшицева градиента, поскольку
\[
    \sup_{x\in\RR}f''(x)=+\infty.
\]
Вблизи минимума градиент \(4x^3\) быстро стремится к нулю, поэтому
градиентный спуск
\[
    x^{k+1}
    =
    x^k-4\alpha(x^k)^3
    =
    x^k\bigl(1-4\alpha(x^k)^2\bigr)
\]
замедляется. При большом шаге или большой начальной точке множитель
\(1-4\alpha(x^k)^2\) может иметь большую абсолютную величину, что
вызывает колебания или расходимость. Поэтому упоминаемую в лекции
граница шага около \(2\) относится только к конкретного эксперимента, а не как универсальную глобальную гарантию.

Метод тяжёлого шара имеет вид
\[
    x^{k+1}
    =
    x^k-4\alpha(x^k)^3
    +\beta(x^k-x^{k-1}),
\]
где \(\alpha>0\) --- шаг, а \(\beta\ge0\) --- коэффициент импульса.
Если
\[
    v^k=x^k-x^{k-1},
\]
то
\[
    v^{k+1}=\beta v^k-4\alpha(x^k)^3,
    \qquad
    x^{k+1}=x^k+v^{k+1}.
\]
В демонстрации используются, в частности,
\[
    x^0=-0.5,
    \qquad
    \alpha=0.5,
    \qquad
    \beta=0.4.
\]
Малое значение \(\beta\), например \(0.1\), почти не отличается от
обычного градиентного шага. Большое значение, например \(0.9\), может
ускорить движение в плоской области, но создаёт выраженные колебания
и немонотонность. Поэтому \(\alpha\) и \(\beta\) должны
согласовываться с гладкостью, кривизной и начальной точкой. Для
квадратичных сильно выпуклых функций параметры можно связать с
крайними собственными значениями гессиана; для общей функции обычно
нужны оценки \(L\) и \(\mu\) либо численная настройка.

Поскольку \(f^\star=0\), величина
\[
    f(x^k)=(x^k)^4
\]
непосредственно измеряет ошибку по значению функции. График
\(k\mapsto\log f(x^k)\) показывает различие в порядках величины и
выявляет временные увеличения ошибки, вызванные чрезмерным импульсом.

\section{Переход к непрерывной динамике}

Из дискретного обновления
\[
    x^{k+1}=x^k-\eta\nabla f(x^k)
\]
получаем
\[
    \frac{x^{k+1}-x^k}{\eta}
    =
    -\nabla f(x^k).
\]
В пределе малого шага возникает градиентный поток.

\begin{definition}[Градиентный поток]
Динамическая система
\[
    \dot x=-\nabla f(x)
\]
называется градиентным потоком функции \(f\).
\end{definition}

Его равновесия удовлетворяют
\[
    \nabla f(x^{\mathrm e})=0.
\]
Для выпуклой функции они являются глобальными минимумами, для строго
выпуклой --- единственным минимумом, а без выпуклости могут быть
локальными минимумами, максимумами или седловыми точками.

Непрерывная постановка позволяет проектировать и другие поля,
например
\[
    \dot x
    =
    -\frac{\nabla f(x)}
    {\lVert\nabla f(x)\rVert+\varepsilon},
    \qquad
    \varepsilon>0.
\]
Формальный выбор \(\varepsilon=0\) даёт единичную скорость вне
стационарных точек. Для
\[
    f(x)=\frac12x^2
\]
получается
\[
    \dot x=-\operatorname{sign}(x),
\]
и нуль достигается за конечное время. При \(\varepsilon>0\) поле
становится регулярнее, но точное конечно-временное достижение обычно
заменяется асимптотическим приближением.

Общая динамическая система записывается как
\[
    \dot x=g(t,x).
\]
Она называется неавтономной, если правая часть явно зависит от
времени, и автономной, если
\[
    \dot x=g(x).
\]
Точка \(x^{\mathrm e}\) является равновесием автономной системы, если
\[
    g(x^{\mathrm e})=0.
\]
Для неавтономной системы постоянное равновесие должно удовлетворять
\(g(t,x^{\mathrm e})=0\) для всех \(t\). После замены
\[
    z=x-x^{\mathrm e}
\]
равновесие переносится в начало координат. Например,
\[
    \dot x=-2(x-1)
\]
переходит в
\[
    \dot z=-2z.
\]

Непрерывные модели удобны для анализа и проектирования, но свойства
их дискретизации требуют отдельной проверки: устойчивое непрерывное
поле может потерять устойчивость при слишком крупном шаге.

\section{Устойчивость, притягательность и скорость}

Рассмотрим систему
\[
    \dot x=g(t,x),
    \qquad
    g(t,0)=0.
\]

\begin{definition}[Устойчивость по Ляпунову]
Равновесие \(x=0\) называется устойчивым, если для каждого
\(\varepsilon>0\) и каждого начального момента \(t_0\) существует
\(\delta(\varepsilon,t_0)>0\) такое, что
\[
    \lVert x(t_0)\rVert<\delta
    \quad\Longrightarrow\quad
    \lVert x(t)\rVert<\varepsilon
\]
для всех \(t\ge t_0\).
\end{definition}

Устойчивость означает сохранение близости, но не требует
\(x(t)\to0\). Она называется равномерной, если \(\delta\) можно
выбрать независимо от \(t_0\); для автономных систем эта независимость
обычно следует из инвариантности относительно сдвига времени.

\begin{definition}[Притягательность]
Равновесие локально притягивает траектории, если существует \(r>0\),
такое что
\[
    \lVert x(t_0)\rVert<r
    \quad\Longrightarrow\quad
    \lim_{t\to\infty}x(t)=0.
\]
Если предел выполняется для всех начальных состояний,
притягательность называется глобальной.
\end{definition}

Притягательность не контролирует промежуточные отклонения. Поэтому
равновесие может притягивать траектории, но не быть устойчивым.

\begin{definition}[Асимптотическая устойчивость]
Равновесие называется локально асимптотически устойчивым, если оно
устойчиво по Ляпунову и локально притягивает траектории. Если
притяжение глобально, говорят о глобальной асимптотической
устойчивости.
\end{definition}

\begin{definition}[Экспоненциальная устойчивость]
Равновесие локально экспоненциально устойчиво, если существуют
\(M\ge1\), \(\alpha>0\) и \(r>0\), для которых
\[
    \lVert x(t)\rVert
    \le
    M\exp\bigl(-\alpha(t-t_0)\bigr)
    \lVert x(t_0)\rVert
\]
при \(\lVert x(t_0)\rVert<r\) и всех \(t\ge t_0\). Если оценка
справедлива для любых начальных состояний, устойчивость глобальная.
\end{definition}

Экспоненциальная устойчивость влечёт асимптотическую и даёт время
достижения точности порядка
\[
    \frac1\alpha\log\frac1\varepsilon.
\]
Простейший пример ---
\[
    \dot x=-ax,
    \qquad
    a>0,
\]
для которого
\[
    x(t)=e^{-a(t-t_0)}x(t_0).
\]
Это одновременно градиентный поток функции
\[
    f(x)=\frac a2x^2.
\]

\section{Конечное и фиксированное время}

\begin{definition}[Конечно-временная сходимость]
Равновесие достигается за конечное время, если для каждого начального
состояния из некоторой области существует
\[
    T(x(t_0))<\infty,
\]
такое что
\[
    x(t)=0
\]
для всех \(t\ge t_0+T(x(t_0))\). Если равновесие при этом устойчиво
по Ляпунову, оно называется конечно-временно устойчивым.
\end{definition}

Для системы
\[
    \dot x=-\operatorname{sign}(x)
\]
имеем
\[
    T(x_0)=|x_0|,
\]
поэтому время зависит от начального состояния.

\begin{definition}[Фиксированно-временная сходимость]
Равновесие достигается за фиксированное время, если существует
\(T_{\max}<\infty\), такое что
\[
    T(x_0)\le T_{\max}
\]
для всех допустимых начальных состояний. При устойчивости по Ляпунову
говорят о фиксированно-временной устойчивости.
\end{definition}

Асимптотическая устойчивость гарантирует только существование
предела. Экспоненциальная устойчивость задаёт скорость затухания.
Конечно-временная сходимость обеспечивает точное достижение за время,
зависящее от \(x_0\). Фиксированно-временная сходимость даёт единую
верхнюю границу для всех запусков. В распределённых системах это
позволяет заранее ограничить момент завершения согласования.

Для градиентного потока соответствия имеют вид
\[
    \text{равновесие}
    \quad\Longleftrightarrow\quad
    \nabla f(x^\star)=0,
\]
\[
    \text{асимптотическая устойчивость}
    \quad\Longleftrightarrow\quad
    x(t)\to x^\star,
\]
\[
    \text{экспоненциальная устойчивость}
    \quad\Longleftrightarrow\quad
    \lVert x(t)-x^\star\rVert
    \le
    M e^{-\alpha(t-t_0)}
    \lVert x(t_0)-x^\star\rVert.
\]
Конечно- и фиксированно-временная устойчивость означают точное
достижение оптимума после конечного момента.

\section{Прямой метод Ляпунова}

Прямой метод исследует устойчивость без явного решения
дифференциального уравнения. Для этого строится скалярная функция
\(V\colon\RR^n\to\RR\).

\begin{definition}[Положительно определённая функция]
Функция \(V\) положительно определена относительно начала координат,
если
\[
    V(0)=0,
    \qquad
    V(x)>0
\]
для всех \(x\ne0\) в рассматриваемой области.
\end{definition}

Вдоль траектории автономной системы \(\dot x=g(x)\)
\[
    \dot V(x)
    =
    \nabla V(x)^{\mathsf T}g(x).
\]
Если \(\dot V\le0\), обобщённая энергия не возрастает; если
\(\dot V<0\) вне равновесия, она строго убывает.

\begin{theorem}[Прямой метод Ляпунова]
Пусть \(x=0\) --- равновесие системы \(\dot x=g(x)\) и в некоторой
его окрестности существует непрерывно дифференцируемая положительно
определённая функция \(V\). Если
\[
    \dot V(x)\le0,
\]
то равновесие устойчиво по Ляпунову. Если
\[
    \dot V(x)<0
\]
для всех \(x\ne0\), то равновесие локально асимптотически устойчиво.
\end{theorem}

Для глобального вывода обычно дополнительно требуется радиальная
неограниченность:
\[
    V(x)\to+\infty
    \qquad
    \text{при }
    \lVert x\rVert\to\infty.
\]

\begin{theorem}[Экспоненциальная оценка]
Пусть существуют \(c_1,c_2,c_3>0\), такие что
\[
    c_1\lVert x\rVert^2
    \le
    V(x)
    \le
    c_2\lVert x\rVert^2
\]
и
\[
    \dot V(x)
    \le
    -c_3\lVert x\rVert^2.
\]
Тогда равновесие экспоненциально устойчиво.
\end{theorem}

Действительно,
\[
    \dot V
    \le
    -\frac{c_3}{c_2}V,
\]
поэтому
\[
    V(x(t))
    \le
    \exp\left(
        -\frac{c_3}{c_2}(t-t_0)
    \right)
    V(x(t_0)).
\]
Двусторонняя оценка на \(V\) переносит это неравенство на норму
состояния. 

Для механической системы естественным кандидатом служит энергия. В
модели «масса--пружина--демпфер» можно взять
\[
    V(x,\dot x)
    =
    \frac12kx^2+
    \frac12m\dot x^2.
\]
Демпфер обеспечивает \(\dot V\le0\). Однако функция Ляпунова не
обязана иметь физический смысл и может быть специально
сконструированной математической величиной.

\section{Функции Ляпунова в оптимизации}

Для градиентного потока
\[
    \dot x=-\nabla f(x)
\]
естественно выбрать
\[
    V(x)=f(x)-f^\star.
\]
Тогда
\[
    \dot V(x)
    =
    \nabla f(x)^{\mathsf T}\dot x
    =
    -\lVert\nabla f(x)\rVert^2
    \le0.
\]
Если \(\nabla f(x)=0\) только в точке \(x^\star\), производная строго
отрицательна вне минимума; при дополнительных условиях регулярности и
компактности это ведёт к асимптотической устойчивости.

Если функция удовлетворяет PL-неравенству
\[
    \frac12\lVert\nabla f(x)\rVert^2
    \ge
    \mu\bigl(f(x)-f^\star\bigr),
\]
то
\[
    \dot V
    \le
    -2\mu V
\]
и
\[
    V(x(t))
    \le
    e^{-2\mu(t-t_0)}V(x(t_0)).
\]
Следовательно, ошибка по значению функции убывает экспоненциально.
Другой стандартный кандидат ---
\[
    V(x)=\frac12\lVert x-x^\star\rVert^2.
\]

Для алгоритмов с импульсом, ограничениями или сетевым взаимодействием
функция Ляпунова может одновременно включать ошибку по значению,
расстояние до решения, скорость, двойственные переменные и ошибку
согласования. Главная трудность прямого метода состоит именно в
выборе такой комбинации. Универсального простого правила, выдающего
удобную аналитическую функцию для любой нелинейной системы, нет.

Вместе с тем существуют обратные теоремы Ляпунова, гарантирующие
существование функции при различных видах устойчивости; сложной
остаётся её явная и удобная для расчётов конструкция.
